Os ciclones extratropicais (ETCs) são sistemas recorrentes na costa do Atlântico Sul-Ocidental que geram ventos e precipitações, cuja coincidência espacial costuma ser assumida.
Enquanto o Hemisfério Norte conta com estudos sistemáticos sobre como esses campos se distribuem, no Hemisfério Sul é mais limitado o uso da densidade espacial e das Funções Ortogonais Empíricas (EOF) univariadas e multivariadas (MEOF) nas diferentes fases do desenvolvimento dos ETCs.
Esta tese desenvolve tal integração para avaliar se a intensidade dinâmica do ETC prediz a magnitude e a localização de seus máximos de vento a 10 metros e precipitação, ou se ambos os campos desenvolvem uma organização estrutural distinta segundo a fase evolutiva.

A análise se constrói sobre o catálogo de trajetórias por vorticidade relativa a 850~hPa \citep{gramcianinov2020analysis} e a classificação objetiva do ciclo de vida em quatro fases mediante \textit{CycloPhaser} \citep{coutodesouza2024new, deSouza2025CycloPhaser}, sobre o ERA5 (1979--2020).
A amostra foi estratificada por três regiões de ciclogênese (SBR, LPB, ARG) e dois estratos de intensidade: a População Global e o subconjunto dos ciclones mais intensos (p90, extraído da População Global).
Sobre essa base foram aplicadas, de forma progressiva, funções de densidade de probabilidade (PDF) de magnitude, estimação de densidade espacial (PDFe) de localização, decomposição EOF univariada com avaliação de significância mediante Bootstrap-$t$, e decomposição MEOF da covariância conjunta, integradas finalmente em protótipos estruturais.

Na População Global, um único modo EOF domina a variabilidade do vento (31--57\% da variância) e se correlaciona com a vorticidade; a precipitação, por sua vez, se distribui entre vários modos (EOF1 de apenas 11--16\%).
Essa organização muda no subconjunto p90. O vento se concentra no quadrante noroeste (moda de 23--25~m\,s$^{-1}$, densidades de até $27\times10^{-4}$ em SBR), enquanto, a partir da maturidade, os máximos de precipitação se afastam do centro e formam uma estrutura em arco a leste e sudeste.
Na População Global, o MEOF1 captura 14.8--26.6\% da covariância conjunta, com os gradientes de precipitação e vento em eixos perpendiculares. Em p90 cai para 11.65--23.47\% e, na maturidade, situa ambos os campos em setores opostos do ciclone, separação respaldada pela correlação significativa entre as PC1 univariadas de ambos os campos ($|r|$ entre $0.31$ e $0.45$), embora o vento domine esse modo conjunto na maioria dos casos.
Os protótipos estruturais quantificam essa separação em mais de $8^{\circ}$--$9^{\circ}$ durante a maturidade e o decaimento, com ARG alcançando a maior magnitude (${\sim}10$--$11^{\circ}$).
A magnitude dessa separação varia regionalmente sem um único fator identificado, embora sua direção, precipitação a leste-sudeste e vento a noroeste, se mantenha consistente nas três regiões.

Esses resultados mostram que, nos ETCs mais intensos do Atlântico Sul-Ocidental, os máximos de vento e precipitação se situam em quadrantes opostos durante a maturidade, e que a vorticidade do centro, que segue associada à organização do vento, deixa de descrever a da precipitação em SBR e LPB.
Esta tese combina PDFe, EOF univariado e MEOF para quantificar essa separação por fase e região de gênese no Atlântico Sul-Ocidental, com implicações para o risco costeiro no Brasil, Uruguai e Argentina.

\vspace{1cm}

\noindent \textbf{Palavras-chave:} América do Sul; ciclones extratropicais; vento a 10 metros; precipitação; estrutura espacial; EOF multivariada.
